% =============================================================================
% Project : THE_PHILOSOPHY_OF_VIABILITY
% Book    : Book I — Ontology of Viability
% Book ID : VS.B01
%
% File    : 2_MAINMATTER/PART_I_IDENTITY_AND_MANIFESTATION/
%           CHAPTER_02_INVARIANCE/sections/
%           section_09_forms_and_domains_of_invariance.tex
%
% Section : 09 — Forms and Domains of Invariance
% Chapter : 02 — Invariance
%
% Purpose :
%   To distinguish the analytical forms through which preservation
%   may be described from the domains in which invariance claims
%   are examined, and to clarify the limits of transferring
%   invariance claims between formal and ontological contexts.
%
% Conceptual Scope :
%   Forms of invariance, domains of invariance, formal domains,
%   ontological domains, four fundamental domains, three
%   intermediate hubs, cross-domain invariance, local and
%   cross-domain transformation, and the scope of invariance claims.
%
% =============================================================================

\section{Forms and Domains of Invariance}
\label{sec:forms-and-domains-of-invariance}

\subsection{The Need to Distinguish Forms from Domains}

The analysis of invariance requires the distinction between
two analytical dimensions: the forms through which preservation
is described and the domains within which a claim of preservation
is examined.

These dimensions address different questions. The forms of
invariance identify the kind of determination that is claimed
to remain preserved. The domains of invariance identify the
ontological or analytical field within which that preservation
is being considered.

The distinction is necessary because the description of what
remains preserved does not, by itself, determine the ontological
domain to which the determination belongs. Conversely, identifying
a domain does not establish which determination is preserved
within it or which form of invariance applies.

For example, structural invariance identifies a claim concerning
the preservation of an organizational determination. It does
not, by itself, establish whether the relevant structure belongs
to a formal system, an existent's structural determination,
or another context in which structural organization is examined.

A rigorous account must therefore specify both dimensions
without allowing one to substitute for the other. It must
also identify the transformation, respect, criterion, and
scope relevant to the particular claim.

\subsection{Forms of Invariance}

Forms of invariance identify analytically distinguishable
ways in which a determination may be claimed to remain
preserved under transformation. They concern the respect
in which preservation is asserted, not separate ontological
entities or necessarily independent kinds of being.

The following forms provide a working analytical vocabulary:

\begin{itemize}
    \item \textbf{Determinational Invariance:}
          The preservation of a specified determination
          or feature across a defined transformation.

    \item \textbf{Structural Invariance:}
          The preservation of a specified organizational
          pattern, arrangement, or structural constraint.

    \item \textbf{Relational Invariance:}
          The preservation of a specified relation between
          relata under a defined transformation.

    \item \textbf{Compositional Invariance:}
          The preservation of a specified determination
          concerning the composition of an existent.

    \item \textbf{Quantitative Invariance:}
          The preservation of a specified quantitative
          determination, such as magnitude, number, or ratio.

    \item \textbf{Qualitative Invariance:}
          The preservation of a specified qualitative
          determination under the transformation considered.

    \item \textbf{Configurational Invariance:}
          The preservation of a specified configuration
          or arrangement of determinations.

    \item \textbf{Relational-Structural Invariance:}
          The preservation of an organized pattern of
          relations, even where particular relata or
          individual relations may change, provided that
          the relevant organizational criterion remains
          satisfied.
\end{itemize}

These forms are analytical distinctions rather than an
exhaustive taxonomy. Their purpose is to clarify the respect
in which a preservation claim is made. They do not imply
that every existent exhibits every form, that every
transformation preserves at least one form, or that each
form has an independent ontological existence.

The classification also does not imply that the forms
are mutually exclusive. A single case may involve more
than one kind of preservation, depending on the determinations
and criteria under examination.

% =============================================================================
% Project : THE_PHILOSOPHY_OF_VIABILITY
% Book    : Book I — Ontology of Viability
% Book ID : VS.B01
%
% File    : 4_ASSETS/tables/
%           VS.B01.P01.C02.S09.SS02.P004.TL01.analytical_forms_of_invariance.tex
%
% Asset   : TL01 — Analytical Forms of Invariance
% Location: Chapter 02 / Section 09 / Subsection 02 / Paragraph 004
%
% Purpose :
%   To distinguish the principal analytical forms of invariance
%   by identifying their respective objects of preservation
%   and limits of inference.
%
% =============================================================================

\begin{table}[htbp]
    \centering
    \small
    \renewcommand{\arraystretch}{1.35}
    \setlength{\tabcolsep}{5pt}

    \begin{tabular}{|p{3.0cm}|p{5.2cm}|p{5.2cm}|}
        \hline
        \textbf{Form} &
        \textbf{Determination Examined} &
        \textbf{Limit of Inference}
        \\
        \hline

        Determinational &
        A specified determination or feature. &
        Does not establish the preservation of the
        existent as a whole or its numerical identity.
        \\
        \hline

        Structural &
        An organizational pattern, arrangement,
        or structural constraint. &
        Does not establish the preservation of
        every component or numerical identity.
        \\
        \hline

        Relational &
        A specified relation between relata. &
        Does not establish that every relatum
        remains unchanged or numerically identical.
        \\
        \hline

        Compositional &
        A specified determination concerning
        the composition of an existent. &
        Does not establish the preservation of
        every component or the whole existent's identity.
        \\
        \hline

        Quantitative &
        A specified magnitude, number, ratio,
        or other quantitative determination. &
        Does not establish qualitative, structural,
        or ontological preservation beyond its criterion.
        \\
        \hline

        Qualitative &
        A specified qualitative determination. &
        Does not establish quantitative or
        numerical identity.
        \\
        \hline

        Configurational &
        A specified configuration or arrangement
        of determinations. &
        Does not establish the preservation of
        every determination or component.
        \\
        \hline

        Relational-Structural &
        An organized pattern of relations under
        a specified organizational criterion. &
        Does not establish that every individual
        relation or relatum remains unchanged.
        \\
        \hline
    \end{tabular}

    \caption{Analytical forms of invariance and the respective
    determinations they examine. The listed forms are not
    mutually exclusive and do not constitute an exhaustive
    ontological taxonomy.}
    \label{tab:analytical-forms-of-invariance}
\end{table}

% =============================================================================
% END OF TABLE TL01
% =============================================================================

% =============================================================================
% Project : THE_PHILOSOPHY_OF_VIABILITY
% Book    : Book I — Ontology of Viability
% Book ID : VS.B01
%
% File    : 4_ASSETS/diagrams/
%           VS.B01.P01.C02.S09.SS02.P004.DG01.forms_and_domains_of_invariance.tex
%
% Asset   : DG01 — Forms and Domains of Invariance
% Location: Chapter 02 / Section 09 / Subsection 02 / Paragraph 004
%
% Purpose :
%   To visualize the analytical independence of forms of
%   invariance and domains of invariance.
%
% Conceptual Caution :
%   The figure represents possible analytical cross-classification.
%   It does not assert that every combination is instantiated
%   or ontologically established.
%
% Requires (preamble):
%   \usepackage{tikz}
%   \usetikzlibrary{calc, positioning, fit}
% =============================================================================

\begin{figure}[htbp]
    \centering

    \begin{tikzpicture}[
        font=\small,
        node distance=0.4cm,
        % --- base box ---
        box/.style={
            draw, rounded corners=3pt, align=center,
            inner sep=5pt, line width=0.6pt
        },
        % --- node roles (names avoid clashes with built-in TikZ keys) ---
        formbox/.style={box, text width=3.4cm, minimum height=2.3cm,
                        fill=blue!8, draw=blue!55!black},
        domainbox/.style={box, text width=2.4cm, minimum height=1.4cm,
                          fill=gray!10, draw=black!60},
        % --- frames, titles, labels ---
        vsgroup/.style={draw=black!45, dashed, rounded corners=5pt,
                        inner sep=8pt, line width=0.5pt},
        axistitle/.style={font=\small\bfseries, align=center},
        grouplabel/.style={font=\footnotesize\itshape, text=black!75,
                           align=center},
        % --- undirected link (no arrowheads: no priority between axes) ---
        vslink/.style={draw=black!50, dashed, line width=0.6pt}
    ]

    % ---------------------------------------------------------------
    % Form axis (left)
    % ---------------------------------------------------------------
    \node[formbox] (question) {
        What determination is claimed to remain preserved?
    };
    \node[formbox, below=of question] (examples) {
        Structural, relational, quantitative, qualitative,
        compositional, and other analytical forms
    };

    % ---------------------------------------------------------------
    % Domain axis (right): two columns of three, top-aligned with Form
    % ---------------------------------------------------------------
    \node[domainbox, anchor=north west] (d1) at
        ($(question.north east)+(1.8,0)$) {Formal Domain};
    \node[domainbox, below=of d1] (d2) {First Determination};
    \node[domainbox, below=of d2] (d3) {Structural Determination};

    \node[domainbox, right=of d1] (d4) {Relational Manifestation};
    \node[domainbox, below=of d4] (d5) {Material Manifestation};
    \node[domainbox, below=of d5] (d6) {Intermediate Hubs};

    % ---------------------------------------------------------------
    % Frames, titles, labels
    % ---------------------------------------------------------------
    \node[vsgroup, fit=(question)(examples)] (groupF) {};
    \node[vsgroup, fit=(d1)(d2)(d3)(d4)(d5)(d6)] (groupD) {};

    \node[axistitle, above=0.2cm of groupF] {FORM OF INVARIANCE};
    \node[axistitle, above=0.2cm of groupD] {DOMAIN OF EXAMINATION};

    \node[grouplabel, below=0.15cm of groupD] {
        Distinct domains of examination
    };

    % ---------------------------------------------------------------
    % Independence marker (symmetric, no arrowheads)
    % ---------------------------------------------------------------
    \node[circle, draw=black!60, fill=white, inner sep=2pt, font=\large]
        (times) at ($(groupF.east)!0.5!(groupD.west)$) {$\times$};
    \draw[vslink] (groupF.east) -- (times);
    \draw[vslink] (times) -- (groupD.west);

    % ---------------------------------------------------------------
    % Note (centred under both axes)
    % ---------------------------------------------------------------
    \coordinate (mid) at ($(groupF.south)!0.5!(groupD.south)$);
    \node[below=0.9cm of mid, text width=10.5cm, align=center,
          font=\footnotesize\itshape] {
        The form identifies the respect of preservation;
        the domain identifies the field in which that
        preservation is examined. Neither dimension
        determines the other.
    };

    \end{tikzpicture}

    \caption{Forms and domains constitute distinct analytical
    dimensions. The diagram presents their independence without
    asserting that every possible form-domain combination has
    been established.}
    \label{fig:forms-and-domains-of-invariance}
\end{figure}

% =============================================================================
% END OF DIAGRAM DG01
% =============================================================================

\subsection{The Non-Exclusivity of Forms}

The analytical forms of invariance may overlap because
a particular determination can be examined under more
than one respect.

A quantitative determination may be involved in a
structural arrangement. A relation may contribute to
a compositional organization. A configuration may
include qualitative and quantitative determinations
at the same time.

Such overlap does not make the distinctions redundant.
It indicates that different analytical descriptions
may address different aspects of one case of preservation.
The descriptions must nevertheless specify their respective
criteria so that the same feature is not counted several
times as though several independent invariants had
necessarily been established.

The forms should therefore not be interpreted as a
rigid classification in which every case must belong
to exactly one category. Nor should their analytical
distinction be taken to establish separate ontological
layers or independent entities.

Their usefulness depends on whether they clarify the
respect in which preservation is asserted. Where a
proposed category does not identify a distinct respect,
its status as a separate form requires reconsideration.

\subsection{Domains of Invariance}

A domain of invariance identifies the ontological or
analytical field within which a claim of preservation
is examined.

The domain does not specify, by itself, what is preserved.
It situates the determination under examination and
constrains the interpretation that may be assigned
to its preservation.

The same analytical form may be considered in different
domains. Structural invariance, for example, may be
examined within a formal system or in relation to the
structural determination of an existent. The use of
the same term does not establish that the two cases
have the same ontological status.

Likewise, a quantitative relation preserved within
a mathematical description does not, merely by being
formally valid, establish that the corresponding
determination is preserved in the existent being
described.

The domain must therefore be identified independently
of the form. Any transfer of a conclusion between
domains requires an account of the relation between
the determinations involved.

\subsection{Formal and Ontological Domains}

A distinction must be maintained between formal
domains and ontological domains.

A formal domain is constituted by the elements,
relations, operations, and rules specified within
a formal system. Mathematical, logical, and other
formally defined systems may provide domains in
which invariance can be demonstrated according
to explicit rules.

An ontological domain concerns determinations
belonging to an existent or to the constitution
of reality, independently of whether those
determinations are represented in a formal system.

Formal invariance may be established when a specified
property, relation, or structure is preserved under
a defined class of transformations within a formal
system. Such a result establishes a relation within
the scope of that system.

It does not automatically establish that an existent
instantiates the relevant formal structure, that
the formal transformation corresponds to an
ontological transformation, or that the property
is preserved in the existent itself.

The transition from a formal result to an ontological
claim therefore requires an additional account of
the correspondence between the formal elements
and the determinations of the existent. This account
must clarify what is represented, how the formal
transformation relates to the transformation under
consideration, and why the formal preservation is
relevant to the ontological claim.

This distinction does not diminish the value of
formal analysis. It specifies the scope of its
conclusions and prevents formal validity from
being treated as an independent proof of
ontological instantiation.

\subsection{Ontological Domains and Their Distinctness}

Within the proposed ontological articulation of
this book, four fundamental domains are distinguished:

\begin{enumerate}
    \item \textbf{First Determination}
    \item \textbf{Structural Determination}
    \item \textbf{Relational Manifestation}
    \item \textbf{Material Manifestation}
\end{enumerate}

These domains identify distinct dimensions of
ontological determination and manifestation.
Their distinction is intended to clarify different
aspects of the ontological articulation, not to
divide reality into isolated regions.

An existent may involve determinations relevant
to more than one domain. A transformation may
also affect determinations in several domains,
whether directly or through relations between them.

Nevertheless, the domains must not be treated
as interchangeable. A determination considered
within one domain cannot be assigned the same
ontological status as a determination considered
within another merely because similar language
is used to describe them.

The proposed distinctions also do not establish
that the four domains are exhaustive, mutually
exclusive in every possible application, or
arranged in a temporal or linear sequence.
Their precise relations remain subject to
the broader ontological analysis.

\subsection{The Four Fundamental Domains}

\subsubsection{First Determination}

First Determination concerns a foundational
dimension of determination within the proposed
ontological articulation.

An invariance claim situated in this domain
would concern the preservation of a determination
properly belonging to this dimension. Such a
claim must identify the relevant determination
and the respect in which it is said to remain
preserved.

It must not be inferred merely from the
preservation of a descriptive feature,
structural arrangement, relational pattern,
or material characteristic.

The precise nature and scope of any proposed
invariant within First Determination require
independent clarification. The identification
of this domain does not establish that a
particular invariant exists within it.

\subsubsection{Structural Determination}

Structural Determination concerns the structural
organization and determinate arrangement of
an existent.

Invariance within this domain may concern
the preservation of a specified organizational
pattern, structural relation, or constraint.

Such a claim must identify the relevant
structural respect and the transformation
under consideration. The preservation of
individual components does not automatically
establish the preservation of their organization.
Conversely, the preservation of an organizational
pattern does not necessarily establish that
every component remains unchanged.

Structural invariance must also remain distinct
from the identity of the existent. The preservation
of a structural determination does not independently
establish numerical identity, and structural
transformation does not independently establish
its loss.

\subsubsection{Relational Manifestation}

Relational Manifestation concerns the manifestation
of an existent through determinate relations.

An invariance claim within this domain may
concern the preservation of a specified relation,
relational determination, or organization of
relations.

The preservation of a relation does not
necessarily establish the preservation of
every relatum involved. Likewise, the alteration
of particular relations does not automatically
establish the disappearance of every relational
determination.

A claim must therefore specify the relata,
the relation under consideration, and the
respect in which preservation is asserted.
Where the organization of relations is involved,
it must be distinguished from the preservation
of each individual relation.

\subsubsection{Material Manifestation}

Material Manifestation concerns the material
manifestation of an existent.

An invariance claim within this domain may
concern specified material determinations,
quantities, configurations, or compositional
aspects.

Material invariance must not be equated with
immutability. A material existent may undergo
transformations while particular material
determinations remain preserved. Conversely,
the preservation of a material determination
does not establish that the existent itself
remains numerically identical.

The claim must identify the material determination,
the transformation, the criterion of preservation,
and the limits within which the preservation
is asserted. Material preservation cannot
be generalized beyond that scope without
additional justification.

\subsection{The Three Intermediate Hubs}

In addition to the four fundamental domains,
the proposed articulation distinguishes three
intermediate hubs:

\begin{enumerate}
    \item \textbf{Hub of Decree}
    \item \textbf{Hub of Signification}
    \item \textbf{Hub of Realization}
\end{enumerate}

These hubs are not presented as three additional
fundamental domains equivalent in status to
the four domains already identified. They
designate intermediate dimensions whose
ontological functions and relations require
their own analysis.

Their inclusion does not establish a linear
sequence connecting all seven designations.
Nor does it establish that each hub occupies
a fixed position between two particular
fundamental domains.

An invariance claim associated with a hub
must be formulated according to the
determination and function being examined.
The hubs must not be collapsed into the
four fundamental domains, and their
distinctive roles must not be erased by
treating all seven designations as equivalent
levels in a single hierarchy.

% =============================================================================
% Project : THE_PHILOSOPHY_OF_VIABILITY
% Book    : Book I — Ontology of Viability
% Book ID : VS.B01
%
% File    : 4_ASSETS/illustrations/
%           VS.B01.P01.C02.S09.SS08.P004.IL01.fundamental_domains_and_intermediate_hubs.tex
%
% Asset   : IL01 — Fundamental Domains and Intermediate Hubs
% Location: Chapter 02 / Section 09 / Subsection 08 / Paragraph 004
%
% Purpose :
%   To provide a visual orientation to the four fundamental
%   domains and three intermediate hubs in the proposed
%   ontological articulation.
%
% Conceptual Caution :
%   Spatial arrangement is illustrative only. It does not
%   establish hierarchy, temporal order, causal direction,
%   exhaustive classification, or fixed relations between
%   the domains and hubs.
%
% Requires (preamble):
%   \usepackage{tikz}
%   \usetikzlibrary{calc, positioning}
% =============================================================================

\begin{figure}[htbp]
    \centering

    % layout parameters (cm)
    \def\dx{4.0}   % horizontal offset of the fundamental domains from centre
    \def\dy{2.0}   % vertical offset of the top / bottom rows from centre
    \def\hx{1.8}   % horizontal offset of the two middle hubs from centre

    \begin{tikzpicture}[
        font=\small,
        node distance=0.4cm,
        % --- node roles ---
        fundamental/.style={
            draw=blue!55!black, fill=blue!8, rounded corners=5pt,
            line width=0.8pt, text width=3.2cm, minimum height=1.2cm,
            align=center, inner sep=6pt
        },
        hub/.style={
            draw=black!60, fill=gray!10, rounded corners=9pt, dashed,
            line width=0.8pt, text width=2.6cm, minimum height=1.2cm,
            align=center, inner sep=6pt
        },
        % --- non-directional association (no arrowheads) ---
        fieldline/.style={draw=black!55, densely dotted, line width=0.8pt},
        % --- legend swatches ---
        legfund/.style={draw=blue!55!black, fill=blue!8, rounded corners=2pt,
                        line width=0.6pt},
        leghub/.style={draw=black!60, fill=gray!10, rounded corners=4pt,
                       dashed, line width=0.6pt}
    ]

    % ---------------------------------------------------------------
    % Fundamental domains (four corners)
    % ---------------------------------------------------------------
    \node[fundamental] (first)      at (-\dx,  \dy) {\textbf{First Determination}};
    \node[fundamental] (structural) at ( \dx,  \dy) {\textbf{Structural Determination}};
    \node[fundamental] (relational) at (-\dx, -\dy) {\textbf{Relational Manifestation}};
    \node[fundamental] (material)   at ( \dx, -\dy) {\textbf{Material Manifestation}};

    % ---------------------------------------------------------------
    % Intermediate hubs (decree aligned with the top row)
    % ---------------------------------------------------------------
    \node[hub] (decree)        at (0,    \dy) {\textbf{Hub of Decree}};
    \node[hub] (signification) at (-\hx, 0)   {\textbf{Hub of Signification}};
    \node[hub] (realization)   at ( \hx, 0)   {\textbf{Hub of Realization}};

    % ---------------------------------------------------------------
    % Non-directional illustrative associations
    % ---------------------------------------------------------------
    \draw[fieldline] (first)      -- (decree);
    \draw[fieldline] (structural) -- (decree);
    \draw[fieldline] (first)      -- (signification);
    \draw[fieldline] (structural) -- (realization);
    \draw[fieldline] (relational) -- (signification);
    \draw[fieldline] (material)   -- (realization);
    \draw[fieldline] (relational) -- (realization);
    \draw[fieldline] (material)   -- (signification);

    % ---------------------------------------------------------------
    % Legend and note (centred under the bottom row)
    % ---------------------------------------------------------------
    \coordinate (mid) at ($(relational.south)!0.5!(material.south)$);

    \node[below=0.8cm of mid, font=\scriptsize] (legend) {
        \tikz[baseline=-0.5ex]\draw[legfund] (0,-0.12) rectangle (0.55,0.12);~Fundamental domain
        \quad\quad
        \tikz[baseline=-0.5ex]\draw[leghub] (0,-0.12) rectangle (0.55,0.12);~Intermediate hub
        \quad\quad
        \tikz[baseline=-0.5ex]\draw[fieldline] (0,0) -- (0.55,0);~Non-directional association
    };

    \node[below=0.45cm of legend, text width=10.5cm, align=center,
          font=\footnotesize\itshape] {
        Four fundamental domains and three intermediate hubs
        within one proposed ontological articulation.
        Their spatial arrangement does not establish a
        linear sequence or a hierarchy.
    };

    \end{tikzpicture}

    \caption{An illustrative orientation to the four fundamental
    domains and three intermediate hubs. The dotted associations
    are non-directional and do not specify fixed relations,
    causal connections, or ontological priority.}
    \label{fig:fundamental-domains-and-intermediate-hubs}
\end{figure}

% =============================================================================
% END OF ILLUSTRATION IL01
% =============================================================================

\subsubsection{Hub of Decree}

Hub of Decree identifies an intermediate
ontological dimension associated with decree.

Any proposed invariance within this hub
must concern a determination properly
belonging to its ontological function.
It cannot be assumed that preservation
within this dimension is equivalent to
the preservation of structural, relational,
or material determinations.

The specific criteria of invariance must
be established according to the nature
of the determination involved. The name
of the hub alone does not establish
which determinations are invariant or
how they are preserved.

\subsubsection{Hub of Signification}

Hub of Signification identifies an intermediate
ontological dimension associated with
signification.

It must not be reduced to language,
subjective meaning, mental representation,
or an act of interpretation. Its ontological
scope cannot be determined solely by
examining linguistic expressions or
psychological processes.

Consequently, invariance within this hub
cannot be defined exclusively through
the persistence of linguistic forms or
mental representations. Any claim must
address the determinations associated
with signification itself, while leaving
open the precise account of their
ontological constitution.

\subsubsection{Hub of Realization}

Hub of Realization identifies an intermediate
ontological dimension associated with
realization.

Invariance within this hub must be examined
according to the particular determinations
involved in realization. It should not
automatically be identified with material
persistence or with the preservation of
an already manifested configuration.

The relation between realization and
the other ontological dimensions requires
independent analysis. The designation
of this hub does not, by itself, establish
a universal mechanism or sequence through
which realization occurs.

\subsection{Cross-Domain Invariance}

An invariance claim may concern determinations
associated with more than one ontological domain.
Such a claim raises a question that must be
distinguished from the mere recurrence of
similar descriptions: are the same determinations
preserved across domains, or are different
determinations related through corresponding
forms of preservation?

These possibilities are not equivalent.

A determination that appears structurally
similar in two domains may have different
ontological functions, conditions, or criteria
of preservation. Similarity between descriptions
does not establish that one and the same
determination is preserved across the domains.

A cross-domain claim therefore requires:

\begin{itemize}
    \item identification of the determination
          considered in each domain;
    \item specification of the transformation
          relevant to each determination;
    \item clarification of the respect in which
          preservation is asserted;
    \item an account of the relation between
          the determinations across domains;
    \item justification for treating the result
          as cross-domain invariance rather than
          analogy or descriptive similarity.
\end{itemize}

Without such clarification, an apparent
cross-domain invariant may result from
similarity in the language of description
rather than an established ontological relation.

Cross-domain invariance is therefore a
claim requiring additional justification,
not a consequence of identifying comparable
features in different domains.

\subsection{Local and Cross-Domain Transformations}

Transformations may be examined within one
specified domain or in relation to determinations
extending across more than one domain.

A local transformation is one whose relevant
scope is restricted to a specified domain.
This does not imply that the transformation
is necessarily without consequences elsewhere;
it means that the claim under examination
is limited to the domain identified.

A cross-domain transformation involves
changes or dependencies that extend across
multiple domains. The relations through
which such changes are connected must
be specified rather than inferred from
their simultaneous occurrence.

This distinction is important because an
invariant may remain preserved within
one domain while relations between that
domain and others are altered.

Likewise, a transformation in one domain
may be associated with consequences in
another without producing an equivalent
transformation there.

Local preservation must therefore not
be confused with global preservation.
A claim of global invariance requires
an account of the relevant determinations
and transformations across every domain
included within its stated scope. It
cannot be inferred from the preservation
of one determination in one domain.

\subsection{The Scope of Domain-Specific Claims}

Every claim of invariance must have a
defined scope. At a minimum, the analysis
must identify the determination under
consideration, the form of invariance
claimed, the domain in which it is examined,
and the transformation and respect relevant
to the claim.

The scope may also require specification
of the criterion of preservation, the
conditions under which the claim applies,
and the limits beyond which the conclusion
is not intended to extend.

An invariant established within a limited
class of transformations cannot automatically
be generalized to all possible transformations.
Similarly, preservation within one domain
does not establish preservation across
the entire existent or across every domain
in which the existent may be considered.

A claim may be valid within a restricted
scope without being universally valid.
Such limitation does not necessarily
invalidate the claim; it determines
what the claim actually establishes.

The scope of validity must therefore
remain explicit throughout the analysis.
Where the scope changes, the conclusion
must be reassessed accordingly.

\subsection{The Limits of a Unified Classification}

The distinction between forms and domains
provides a framework for organizing
invariance claims. It does not imply
that every possible invariant can be
classified exhaustively within a closed
taxonomy.

Further distinctions among forms may
be required as additional modes of
preservation are examined. Likewise,
the analysis of domains may require
further clarification as their ontological
relations become more precise.

The forms listed in this section should
therefore be treated as an analytical
vocabulary rather than a completed
inventory. The four fundamental domains
and three intermediate hubs likewise
belong to a developing ontological
articulation, not an assertion that
every possible domain has already
been identified.

The framework is systematic in its
distinctions but remains open in its
scope. Its purpose is to prevent
conceptual confusion and unjustified
transfers, not to close the inquiry
into every possible form or domain
of invariance.

\subsection{Synthesis}

The analysis establishes two independent
but related dimensions for examining
invariance.

Forms of invariance identify the respect
or kind of determination claimed to
remain preserved. Domains of invariance
identify the ontological or analytical
field within which that preservation
is examined.

The distinction yields the following
conclusions:

\begin{itemize}
    \item A form of invariance does not
          determine its domain.

    \item A domain does not determine
          which form of invariance must
          occur within it.

    \item Formal invariance does not
          automatically establish
          ontological invariance.

    \item Local invariance does not
          establish global invariance.

    \item Invariance within one domain
          does not automatically transfer
          to another domain.

    \item The four fundamental domains
          and three intermediate hubs
          must not be treated as seven
          equivalent or linearly ordered
          ontological levels.

    \item No classification of invariance
          independently establishes the
          identity of an existent.
\end{itemize}

The analysis of invariance must consequently
preserve the distinction between what
is claimed to remain preserved and the
domain in which that preservation is
asserted. It must also distinguish the
formal validity of a preservation claim
from its ontological applicability.

Only through these distinctions can
invariance claims be evaluated without
confusing analytical form, domain,
formal representation, and ontological
determination. The classification
remains a means of clarifying the inquiry,
not a substitute for the arguments
required to establish particular
invariants.

% =============================================================================
% END OF SECTION 09 OF CHAPTER 02
% =============================================================================